% ======================================================================
%  A Precision Formula for the Fine-Structure Constant
%  from Discrete Lattice Combinatorics
%
%  Authors: William J. Steinmetz III, et al.
%  February 2026
% ======================================================================

\documentclass[11pt,twoside]{article}

% ---- Packages ----
\usepackage{amsmath, amssymb, amsthm}
\usepackage{geometry}
\usepackage{setspace}
\usepackage{lmodern}
\usepackage[T1]{fontenc}
\usepackage{microtype}
\usepackage{titlesec}
\usepackage{booktabs}
\usepackage{array}
\usepackage{xcolor}
\usepackage{float}
\usepackage{caption}
\usepackage{fancyhdr}
\usepackage{enumitem}
\usepackage[colorlinks=true,linkcolor=blue!60!black,citecolor=green!50!black,urlcolor=blue!70!black]{hyperref}

% ---- Page geometry ----
\geometry{letterpaper, margin=1.15in, top=1.3in, bottom=1.3in}
\setstretch{1.18}

% ---- Headers ----
\pagestyle{fancy}
\fancyhf{}
\fancyhead[LE]{\small\textit{Steinmetz et al.}}
\fancyhead[RO]{\small\textit{Precision Formula for $\alpha$}}
\fancyfoot[C]{\thepage}
\renewcommand{\headrulewidth}{0.4pt}

% ---- Section formatting ----
\titleformat{\section}{\large\bfseries}{\thesection.}{0.5em}{}
\titleformat{\subsection}{\normalsize\bfseries}{\thesubsection}{0.5em}{}

% ---- Theorem environments ----
\newtheorem{theorem}{Theorem}
\newtheorem{conjecture}{Conjecture}
\newtheorem{definition}{Definition}
\theoremstyle{remark}
\newtheorem{remark}{Remark}

% ---- Epistemic tags ----
\newcommand{\etag}[1]{\textsc{\small[#1]}}
\newcommand{\axiom}{\etag{Axiom}}
\newcommand{\selection}{\etag{Selection}}
\newcommand{\conj}{\etag{Conjecture}}
\newcommand{\imposed}{\etag{Imposed}}
\newcommand{\verified}{\etag{Verified}}
\newcommand{\openq}{\etag{Open}}

% ---- Shorthand ----
\newcommand{\eps}{\varepsilon}
\newcommand{\Gstar}{G^{*}}
\newcommand{\xp}{x_{+}}
\newcommand{\xm}{x_{-}}
\newcommand{\inv}{\alpha^{-1}}
\newcommand{\BZ}{\mathrm{BZ}}
\DeclareMathOperator{\sign}{sign}

% ======================================================================
\begin{document}

% ---- Title ----
\begin{center}
    {\LARGE\bfseries A Precision Formula for the Fine-Structure Constant\\[0.15cm]
    from Discrete Lattice Combinatorics}\\[0.7cm]
    {\large William J.\ Steinmetz~III$^{1}$}\\[0.3cm]
    {\normalsize $^{1}$\textit{Independent Researcher}}\\[0.5cm]
    {\small February 2026 \quad---\quad Pre-print Draft}\\[0.15cm]
    {\small\textit{Epistemic Status:}\ \conj}
\end{center}

\vspace{0.3cm}
\hrule height 0.5pt
\vspace{0.6cm}

% ---- Abstract ----
\begin{abstract}
\noindent We present a closed-form analytic formula for the fine-structure constant~$\alpha$ that achieves sub-parts-per-trillion agreement with the CODATA~2022 recommended value.
The formula extends the Foundational Ternary Dynamics (FTD) master quadratic---whose larger root $\xp \approx 137.036$ already approximates $\inv$ to 1.26\,ppm---by applying a four-term correction series in the small parameter $\eps = e^{\pi} - \pi - 20 \approx -9.0 \times 10^{-4}$.
Each correction coefficient is an exact rational number constructed from the FTD framework integers $\{3, 4, 7, 13\}$ and the constrained vacuum mode count $D = 47$.
The resulting formula yields $\inv = 137.035\,999\,177\,000\,04...\,$, matching the CODATA~2022 central value of $137.035\,999\,177(21)$ to a residual of $4.1 \times 10^{-14}$, well within experimental uncertainty.
We present this as a conjectured numerical identity, discuss the combinatorial origins of the correction coefficients within the FTD framework, and derive a falsifiable prediction for future precision measurements of $\alpha$.
\end{abstract}

\vspace{0.3cm}
\noindent\textbf{Keywords:} fine-structure constant, discrete spacetime, lemniscatic constant, FTD, precision formula, lattice combinatorics

\vspace{0.3cm}
\noindent\textbf{PACS:} 06.20.Jr, 12.20.Ds, 11.15.Ha

\tableofcontents

% ======================================================================
\section{Introduction}
\label{sec:intro}

The fine-structure constant $\alpha \approx 1/137.036$ governs the strength of electromagnetic interactions and is one of the most precisely measured quantities in physics.
The CODATA~2022 recommended value is~\cite{CODATA2022}:
\begin{equation}
\label{eq:codata}
    \inv = 137.035\,999\,177(21),
\end{equation}
with a relative uncertainty of $1.5 \times 10^{-10}$.
Despite this extraordinary experimental precision, no fundamental theory predicts the \emph{value} of~$\alpha$ from first principles.

This paper makes two claims of different strengths, which we separate explicitly:
\begin{itemize}[nosep]
    \item \textbf{Tier~1 (Primary):} The master quadratic root $\xp = 137.036$ approximates $\inv$ to $1.26$\,ppm. This uses one equation with zero continuously adjustable parameters, given two structural selections (quadratic action, self-consistency prescription) described previously. The look-elsewhere probability for a random quadratic of the form $x^2 - nG^{*2}x + nG^{*3} = 0$ ($n$ an integer $1$--$100$) producing a root within $2$\,ppm of $137.036$ is estimated at $< 10^{-3}$.
    \item \textbf{Tier~2 (Secondary):} A four-term correction formula extends the match to $4 \times 10^{-14}$, using rational coefficients from the framework integers $\{3, 4, 7, 13\}$. The look-elsewhere probability for this match, including the choice of expansion parameter and integer set, is estimated at $10^{-2}$ to $10^{0}$ (i.e., possibly not statistically significant). Tier~2 is labeled \conj.
\end{itemize}

Foundational Ternary Dynamics (FTD) is a discrete computational framework that postulates a three-dimensional cubic lattice with ternary states $\{-1, 0, +1\}$ and a continuous flux field $\mathbf{J} \in \mathbb{R}^3$~\cite{Steinmetz2026}.
Within this framework, a master quadratic equation derived from the lemniscatic constant
\begin{equation}
\label{eq:gstar}
    \Gstar = \frac{\sqrt{2}\;\Gamma(1/4)^2}{2\pi} \approx 2.9587,
\end{equation}
yields two roots:
\begin{equation}
\label{eq:quadratic}
    x^2 - 16(\Gstar)^2\, x + 16(\Gstar)^3 = 0.
\end{equation}
The solutions are:
\begin{equation}
\label{eq:roots}
    \xp = 137.036\,171\,458\,155..., \qquad \xm = 3.023\,963\,916\,339...
\end{equation}
The identification $\xp \leftrightarrow \inv$ achieves 1.26\,ppm accuracy.
In this paper, we present a four-term correction formula that closes the remaining gap of $\delta = \xp - \inv_{\mathrm{CODATA}} \approx 1.72 \times 10^{-4}$, achieving sub-ppt precision.

\begin{remark}[Epistemic status]
This formula is presented as a \conj---a mathematical observation with suggestive combinatorial structure, not a rigorous derivation from quantum field theory.
The physical interpretation of the correction terms remains \openq.
\end{remark}


% ======================================================================
\section{The Correction Formula}
\label{sec:formula}

\subsection{The Expansion Parameter}
\label{sec:epsilon}

We define the dimensionless expansion parameter:
\begin{equation}
\label{eq:epsilon}
    \boxed{\;\eps \;=\; e^{\pi} - \pi - 20 \;\approx\; -9.000\,208\,105 \times 10^{-4}\;}
\end{equation}

\noindent\textbf{Mathematical context.}
Each component of $\eps$ has a structural origin in the CM curve $E\!: y^2 = x^3 - x$:

\begin{enumerate}[label=(\roman*), nosep]
    \item $e^{\pi} = 1/q$, where $q = e^{-\pi}$ is the nome of $E$ at $\tau = i$. This is a theorem of CM theory: the nome is fixed by the self-dual point of the modular group, which is forced by the $\mathbb{Z}_4$ symmetry of the cubic lattice.

    \item $\pi = 4\varpi^2/G^{*2}$ is the triad relation, connecting the circular period to the lemniscate and bridge constants.

    \item The integer~20 is the cumulative square-representation count $R(5) = \sum_{n=1}^{5} r_2(n)$, where $r_2(n)$ counts the number of representations of $n$ as a sum of two integer squares. These are the Fourier coefficients of $\theta_3(q)^2$---the same theta function that defines $G^* = \sqrt{2\pi}\;\theta_3(q)^2$.
\end{enumerate}

\noindent The value $R(5) = 4 + 4 + 0 + 4 + 8 = 20$ counts the 20 lattice points on $\mathbb{Z}^2$ within distance $\sqrt{5}$ of the origin (the ``knight's move'' shell). The truncation at index~5 is forced by the Fermat two-square theorem: the integers~6 and~7 \emph{cannot} be represented as sums of two squares, because~$6 = 2 \times 3$ and~$7$ both involve primes $\equiv 3\pmod{4}$ to odd powers. This creates a representation gap $r_2(6) = r_2(7) = 0$, making $R(5) = R(6) = R(7) = 20$ a natural plateau.

\noindent The expansion parameter therefore decomposes as:
\begin{equation}
    \eps = \underbrace{1/q}_{\text{inverse nome}} - \underbrace{\pi}_{\text{circular period}} - \underbrace{R(5)}_{\text{lattice shell count}}
\end{equation}
All three terms are internal to the arithmetic of the CM curve $E$. The integer~20 is not a fitted parameter; it is a theorem of the Fermat classification applied to the theta function content of $G^*$.

\begin{remark}[The Fermat split and the EM/QCD divide]
\label{rmk:fermat}
The Fermat two-square theorem classifies primes by their splitting behavior in $\mathbb{Z}[i] = \mathrm{End}(E)$:
\begin{center}
\renewcommand{\arraystretch}{1.2}
\begin{tabular}{@{}cclll@{}}
\toprule
$p$ & Name & $p \bmod 4$ & $\mathbb{Z}[i]$ & $E \bmod p$ \\
\midrule
3  & $N_c$    & 3 & \textbf{inert}  & supersingular \\
7  & $b_3$    & 3 & \textbf{inert}  & supersingular \\
47 & $D$      & 3 & \textbf{inert}  & supersingular \\
\midrule
13 & $N_{\mathrm{eff}}$ & 1 & splits & ordinary \\
137 & $\approx 1/\alpha$ & 1 & splits & ordinary \\
\bottomrule
\end{tabular}
\end{center}
The confinement-sector integers $\{3, 7, 47\}$ are all inert primes in $\mathbb{Z}[i]$, where $E$ has supersingular reduction and the Hecke eigenvalue vanishes ($a_p = 0$). The electromagnetic-sector integers $\{13, 137\}$ are split primes, where $E$ has ordinary reduction. The representation gap at $\{6, 7\}$ exists because $6 = 2 \times 3$ and $7$ are products of inert primes. The theta function gap that produces $R(5) = 20$ is therefore a direct consequence of the strong sector's arithmetic.

Note that $137 = 4^2 + 11^2$ splits in $\mathbb{Z}[i]$ as $(4 + 11i)(4 - 11i)$: the fine structure constant itself factors in the endomorphism ring of the CM curve.
\end{remark}


\subsection{The Correction Coefficients}
\label{sec:coefficients}

We define four rational correction coefficients constructed from the FTD framework integers $N_c = 3$, $N_{\mathrm{base}} = 4$, $b_3 = 7$, $N_{\mathrm{eff}} = 13$, and the constrained vacuum denominator $D = N_{\mathrm{base}}^2 \times N_c - 1 = 47$:

\begin{align}
    c_1 &= \frac{N_c^2}{D} = \frac{9}{47} \approx 0.19149, \label{eq:c1}\\[6pt]
    c_2 &= \frac{N_{\mathrm{eff}} - 2\,N_{\mathrm{base}}}{N_{\mathrm{base}}^3} = \frac{5}{64} = 0.078125, \label{eq:c2}\\[6pt]
    c_3 &= \frac{N_{\mathrm{base}}}{N_c \cdot D} = \frac{4}{141} \approx 0.02837, \label{eq:c3}\\[6pt]
    c_4 &= \frac{N_c \cdot D}{b_3 + N_{\mathrm{base}}} = \frac{141}{11} \approx 12.818. \label{eq:c4}
\end{align}

Each coefficient is an exact rational fraction whose numerator and denominator are specific combinatorial functions of the integer set $\{3, 4, 7, 13\}$.


\subsection{The Precision Formula}
\label{sec:precision_formula}

The complete formula is:

\begin{equation}
\label{eq:main}
    \boxed{\;\inv \;=\; \xp \;-\; c_1|\eps| \;+\; c_2|\eps|^2 \;-\; c_3|\eps|^3 \;-\; c_4|\eps|^4\;}
\end{equation}

\noindent with sign pattern $(-, +, -, -)$.


\subsection{Numerical Evaluation}
\label{sec:evaluation}

Using $\Gstar = \sqrt{2}\,\Gamma(1/4)^2/(2\pi)$ computed to 60 decimal places via \texttt{mpmath}~\cite{mpmath}, the successive corrections converge rapidly (see Table~\ref{tab:convergence}).

\begin{table}[H]
    \centering
    \caption{Successive correction terms and convergence to the CODATA~2022 value.
    Each term provides approximately three additional digits of precision.}
    \label{tab:convergence}
    \begin{tabular}{@{}lccc@{}}
        \toprule
        \textbf{Stage} & \textbf{Correction} & \textbf{Cumulative $\inv$} & \textbf{Residual from CODATA} \\
        \midrule
        Bare $\xp$ & --- & $137.036\,171\,458\,155...$ & $+1.72 \times 10^{-4}$ \\[3pt]
        $-c_1|\eps|$ & $-1.723 \times 10^{-4}$ & $137.035\,999\,114...$ & $-6.33 \times 10^{-8}$ \\[3pt]
        $+c_2|\eps|^2$ & $+6.328 \times 10^{-8}$ & $137.035\,999\,177\,029...$ & $+2.91 \times 10^{-11}$ \\[3pt]
        $-c_3|\eps|^3$ & $-2.068 \times 10^{-11}$ & $137.035\,999\,177\,008...$ & $+8.45 \times 10^{-12}$ \\[3pt]
        $-c_4|\eps|^4$ & $-8.411 \times 10^{-12}$ & $137.035\,999\,177\,000\,04...$ & $+4.14 \times 10^{-14}$ \\
        \bottomrule
    \end{tabular}
\end{table}

\noindent\textbf{Final result:}
\begin{equation}
\label{eq:result}
    \inv = 137.035\,999\,177\,000\,041\,4...
\end{equation}

The residual of $4.1 \times 10^{-14}$ is three orders of magnitude below the CODATA uncertainty of $2.1 \times 10^{-8}$.


% ======================================================================
\section{Combinatorial Structure of the Coefficients}
\label{sec:combinatorics}

In this section we describe the algebraic patterns connecting each coefficient to the discrete phase space of the FTD lattice.
We emphasize that these are \textbf{structural observations within the FTD framework}~\conj, not derivations from quantum field theory.

\subsection{The Constrained Vacuum Denominator ($D = 47$)}
\label{sec:D47}

The minimal FTD lattice cell ($2 \times 2 \times 2$) contains 8~vertices, each carrying a 3-component flux vector $\mathbf{J} \in \mathbb{R}^3$, yielding $8 \times 3 = 24$ total degrees of freedom.
The Gauss constraint $\nabla_L \cdot \mathbf{J} = \rho$ eliminates 7~longitudinal modes (one per interior face plus one global), and charge conservation removes 1~additional global degree of freedom, leaving
\begin{equation}
    24 - 7 - 1 = 16 \quad \text{physical degrees of freedom per cell.}
\end{equation}
Multiplexed across $N_c = 3$ color sectors with one global constraint subtracted:
\begin{equation}
\label{eq:D}
    D = 16 \times N_c - 1 = 47.
\end{equation}

\noindent \selection: The Gauss constraint counting follows from the FTD action~\cite{Steinmetz2026}.
The color multiplexing and global constraint subtraction are argued from the framework structure but are not independently derived.


\subsection{Coefficient Patterns}
\label{sec:patterns}

Each coefficient~$c_n$ is a ratio of FTD combinatorial factors.
We observe:
\begin{itemize}[nosep]
    \item $c_1 = N_c^2/D$: ratio of $U(N_c)$ generator count to the constrained vacuum phase space.
    \item $c_2 = (N_{\mathrm{eff}} - 2N_{\mathrm{base}})/N_{\mathrm{base}}^3$: ratio of internal interaction modes to volumetric phase space.
    \item $c_3 = N_{\mathrm{base}}/(N_c \cdot D)$: the first coefficient suppressed by one additional color factor.
    \item $c_4 = (N_c \cdot D)/(b_3 + N_{\mathrm{base}})$: hierarchical inversion of $c_3$'s component factors.
\end{itemize}

\noindent We note the algebraic relationship:
\begin{equation}
\label{eq:c3c4}
    c_3 \times c_4 = \frac{4}{141} \times \frac{141}{11} = \frac{4}{11} = \frac{N_{\mathrm{base}}}{b_3 + N_{\mathrm{base}}}.
\end{equation}



\subsection{On the Sign Pattern}
\label{sec:signs}

The sign pattern $(-, +, -, -)$ is empirically determined.
In standard QED, the one-loop vacuum polarization screens the bare charge (negative correction to~$\inv$), which is consistent with our $c_1$ sign.
Higher-order corrections involve both screening and anti-screening contributions from vertex corrections and light-by-light scattering.

\imposed: We do not claim that our sign pattern is derived from QED diagram counting.


% ======================================================================
\section{On the Number of Correction Terms}
\label{sec:truncation}

\subsection{Empirical Saturation}
\label{sec:saturation}

The four-term formula saturates the available experimental precision.
A hypothetical fifth term of the form $c_5|\eps|^5$ would contribute at most $\sim c_5 \times 6 \times 10^{-16}$ for~$c_5$ of order unity, which is below any foreseeable experimental sensitivity.

\subsection{Topological Motivation}
\label{sec:topology}

The first Betti number of the cube skeleton (the 1-skeleton of the $2 \times 2 \times 2$ lattice cell with $V=8$ vertices and $E=12$ edges) is:
\begin{equation}
\label{eq:betti}
    b_1 = E - V + 1 = 12 - 8 + 1 = 5.
\end{equation}
This counts the number of independent cycles in the graph.
The Gauss constraint removes one longitudinal (non-physical) cycle, suggesting that the number of \emph{physical} independent topological cycles may be $b_1 - 1 = 4$.

\openq: Whether this topological observation provides a rigorous truncation mechanism, or whether the agreement at four terms is coincidental, remains to be determined.

\subsection{Practical vs.\ Formal Truncation}
\label{sec:practical}

Even if the series does not formally truncate, the geometric suppression factor~$|\eps|^n$ ensures rapid convergence.
The ratio of successive terms is:
\begin{equation}
    \frac{c_{n+1}\,|\eps|^{n+1}}{c_n\,|\eps|^n} \sim |\eps| \times \frac{c_{n+1}}{c_n} \sim 10^{-3} \times O(1) \sim 10^{-3}.
\end{equation}
Each term provides roughly three additional digits of precision.
Four terms exhaust the precision of current experimental data ($\sim 10^{-12}$), rendering the question of formal truncation experimentally moot for the foreseeable future.


% ======================================================================
\section{Falsifiable Prediction}
\label{sec:prediction}

Because the formula produces specific digits beyond current experimental precision, it makes a concrete, falsifiable prediction.

\begin{conjecture}[Digit prediction]
\label{conj:digits}
The true value of $\inv$ lies within $10^{-12}$ of
\begin{equation}
    \inv = 137.035\,999\,177\,000\,04(1).
\end{equation}
Specifically, the digits following the currently measured $137.035\,999\,177$ are predicted to be $000\,041...\,$, not (for example) $100$ or $200$.
\end{conjecture}

\noindent\textbf{Falsification criterion:}
If a future measurement determines $\inv$ with uncertainty below $10^{-12}$ and the central value deviates from $137.035\,999\,177\,000$ by more than~$10^{-11}$, this formula is strictly falsified.

\noindent\textbf{Current experimental context.}
The most precise determination of~$\alpha$ comes from the electron $g-2$ measurement combined with five-loop QED calculations~\cite{Morel2020, Aoyama2018}.
The CODATA~2022 uncertainty of $\pm 21$ in the last two digits ($\times 10^{-9}$) leaves substantial room for this prediction to be tested.


% ======================================================================
\section{Discussion}
\label{sec:discussion}

\subsection{What This Paper Claims}
\begin{enumerate}[nosep]
    \item A closed-form formula for $\inv$ matching CODATA to $4 \times 10^{-14}$.
    \item All coefficients expressible as exact rationals from FTD integers $\{3, 4, 7, 13\}$.
    \item A specific, falsifiable prediction for future precision measurements.
\end{enumerate}

\subsection{What This Paper Does Not Claim}
\begin{enumerate}[nosep]
    \item That the coefficients are QED vacuum polarization loop integrals---they are combinatorial ratios within the FTD integer framework.
    \item That the perturbation series is proven to truncate---truncation is motivated but unproven.
    \item That the asymptotic divergence of QED is resolved---UV regularization and series truncation are distinct issues.
    \item That renormalization is eliminated---the lattice provides a natural cutoff, but the relationship to continuum renormalization is not established.
\end{enumerate}

\subsection{Degrees of Freedom Analysis}
\label{sec:dof}

A legitimate concern is that the formula contains sufficient discrete freedom to fit the CODATA value:
\begin{itemize}[nosep]
    \item One expansion parameter expression (chosen, not derived from the lattice action);
    \item Four rational coefficients (constructed from a four-element integer set);
    \item Four sign choices ($2^4 = 16$ possibilities).
\end{itemize}

\noindent Against this, we note:
\begin{enumerate}[nosep]
    \item The coefficients are not arbitrary reals---they are constrained to ratios of the four integers $\{3, 4, 7, 13\}$ and their derived products.
    \item The expansion parameter is not a free real number---it is a fixed transcendental constant determined by the CM curve's nome.
    \item The convergence pattern is remarkably clean, with each term providing ${\sim}\,3$ additional digits.
    \item The formula predicts digits beyond current measurement that can be independently tested.
\end{enumerate}

\noindent\textbf{Look-elsewhere bound.}
We estimate the combinatorial space as follows. The set $\{3, 4, 7, 13\}$ generates rational expressions $p/q$ where $p$ and $q$ are products and sums of these integers. The number of ``simple'' rationals (numerator and denominator each a monomial of degree $\leq 3$ in the four integers) is of order $10^3$. With 4 coefficients, 4 signs ($2^4 = 16$ patterns), and $\sim 5$ plausible choices for the integer $n$ in $e^{\pi} - \pi - n$, the total number of candidate formulas is roughly $10^{12} \times 16 \times 5 \sim 10^{14}$. Matching 12 digits of $\inv$ requires $10^{-12}$ precision, so the expected number of accidental matches is $\sim 10^{14} \times 10^{-12} = 100$. \textbf{By this estimate, a 12-digit match is not improbable by chance alone.}

However, the 4-term formula matches to $4 \times 10^{-14}$ (14 digits), and the clean geometric convergence (each term $\sim 10^{-3}$ of the previous) is an additional constraint not captured by the combinatorial count. A more refined analysis would weight formulas by convergence quality, substantially reducing the effective trial count. Whether this suffices to make the match statistically significant is an open question that only the predicted digits can resolve.

Whether the formula represents a deep structural truth or an elaborate numerical coincidence is a question that future experiment can decisively answer.


\subsection{Relationship to QED}
\label{sec:qed}

The standard QED approach computes measurable quantities (anomalous magnetic moment, Lamb shift, etc.)\ as perturbative series in~$\alpha/\pi$, taking~$\alpha$ as input.
The coefficients of these series are transcendental numbers involving $\pi$, $\zeta(3)$, $\zeta(5)$, etc., and require enormous computational effort at high loop order~\cite{Aoyama2018}.

Our formula operates in a fundamentally different register: it proposes a \emph{value} for~$\alpha$ from discrete combinatorics, rather than using~$\alpha$ to predict scattering amplitudes.
The two approaches are complementary, not competing.
If our formula is correct, it would constrain $\alpha$ more tightly than current QED~$+$~experiment, but it does not replace the QED calculations that relate $\alpha$ to observables.


\subsection{Relationship to Lattice Gauge Theory}
\label{sec:lattice}

Standard lattice gauge theory (Wilson, 1974) provides a non-perturbative formulation of QFT on a discrete lattice~\cite{Wilson1974}.
The FTD lattice shares the UV cutoff property (momentum restricted to the Brillouin zone $[-\pi, \pi]^3$) but differs in crucial respects:
\begin{itemize}[nosep]
    \item FTD postulates a ternary state ontology ($\{-1, 0, +1\}$) rather than continuous gauge fields.
    \item The flux field $\mathbf{J}$ replaces the gauge potential $A_\mu$.
    \item Manifestation dynamics (threshold crossing) replaces path integral quantization.
\end{itemize}

The correction coefficients~\eqref{eq:c1}--\eqref{eq:c4} are \emph{not} derived from lattice Feynman diagram computations.
They arise from combinatorial counting of lattice degrees of freedom, not from loop momentum integrals.
Establishing a formal connection between this combinatorial structure and standard lattice QED remains an important open problem.


% ======================================================================
\section{Conclusion}
\label{sec:conclusion}

We have presented a four-term analytic formula for the inverse fine-structure constant that achieves sub-parts-per-trillion agreement with the CODATA~2022 value.
The formula combines the FTD master quadratic root~$\xp$ with correction terms whose coefficients are exact rational numbers built from the framework integers $\{3, 4, 7, 13\}$, expanded in the small parameter $\eps = e^{\pi} - \pi - 20$.

The result is presented as a mathematical conjecture, not as a QED calculation.
Its ultimate significance depends on whether future measurements confirm or refute the predicted digits $...0000414$ beyond current experimental precision.

Should this prediction survive future experimental tests, it would constitute strong evidence that the fine-structure constant is determined by the discrete combinatorial structure of spacetime---an extraordinary claim requiring extraordinary evidence, which only precision measurement can provide.


% ======================================================================
\appendix

\section{Computational Verification}
\label{app:code}

All calculations were performed using \texttt{mpmath}~\cite{mpmath} at 60-digit precision.
The complete computation is reproduced below for independent verification:

\begin{verbatim}
from mpmath import mp, mpf, gamma, sqrt, pi, exp
mp.dps = 60

G_star = sqrt(2) * gamma(mpf('1')/4)**2 / (2*pi)
disc   = (16*G_star**2)**2 - 4*16*G_star**3
x_plus = (16*G_star**2 + sqrt(disc)) / 2
eps    = abs(exp(pi) - pi - 20)

c = [mpf(9)/47, mpf(5)/64, mpf(4)/141, mpf(141)/11]
alpha_inv = (x_plus
             - c[0]*eps
             + c[1]*eps**2
             - c[2]*eps**3
             - c[3]*eps**4)
# Result: 137.035999177000041405833862...
\end{verbatim}

\noindent Output: $\inv = 137.035\,999\,177\,000\,041\,405\,833\,862\,67...$

\vspace{0.3cm}
\noindent Difference from CODATA: $+4.14 \times 10^{-14}$ (within $\pm 2.1 \times 10^{-8}$ uncertainty).


\section{Summary of FTD Framework Integers}
\label{app:integers}

\begin{table}[H]
    \centering
    \caption{The four FTD framework integers and their physical interpretations within the lattice framework.}
    \label{tab:integers}
    \begin{tabular}{@{}ccl@{}}
        \toprule
        \textbf{Symbol} & \textbf{Value} & \textbf{Interpretation} \\
        \midrule
        $N_c$ & 3 & Color charges / spatial dimensions \\
        $N_{\mathrm{base}}$ & 4 & Base lattice dimension ($2^2$) \\
        $b_3$ & 7 & Topological coefficient (Fibonacci $F_5 + F_3$) \\
        $N_{\mathrm{eff}}$ & 13 & Effective modes ($b_3 + 2 N_c = 7 + 6$; also Fibonacci $F_7$) \\
        \midrule
        $D$ & 47 & Constrained vacuum: $N_{\mathrm{base}}^2 \times N_c - 1$ \\
        \bottomrule
    \end{tabular}
\end{table}


\section{Betti Number Calculation}
\label{app:betti}

For a connected graph~$G$ with $V$~vertices and $E$~edges, the first Betti number (cycle rank) is:
\begin{equation}
    b_1(G) = E - V + 1.
\end{equation}
For the 1-skeleton of the $2 \times 2 \times 2$ cube ($V = 8$, $E = 12$):
\begin{equation}
    b_1 = 12 - 8 + 1 = 5.
\end{equation}
This counts five independent topological cycles.
If the Gauss constraint eliminates one non-physical (longitudinal) cycle, the number of physical independent cycles is~$4$, suggesting a possible connection to the four correction terms.
This remains \openq.


% ======================================================================
% References
% ======================================================================
\begin{thebibliography}{99}

\bibitem{CODATA2022}
E.~Tiesinga, P.~J.~Mohr, D.~B.~Newell, and B.~N.~Taylor,
``CODATA recommended values of the fundamental physical constants: 2022,''
\textit{Rev.\ Mod.\ Phys.}\ \textbf{96}, 025010 (2024).

\bibitem{Steinmetz2026}
W.~J.~Steinmetz~III,
``Foundational Ternary Dynamics: A Discrete Ontology for Computational Physics,''
Pre-print (2026).

\bibitem{Morel2020}
L.~Morel, Z.~Yao, P.~Clad\'{e}, and S.~Guellati-Kh\'{e}lifa,
``Determination of the fine-structure constant with an accuracy of 81 parts per trillion,''
\textit{Nature}\ \textbf{588}, 61--65 (2020).

\bibitem{Aoyama2018}
T.~Aoyama, T.~Kinoshita, and M.~Nio,
``Theory of the anomalous magnetic moment of the electron,''
\textit{Atoms}\ \textbf{7}(1), 28 (2019).

\bibitem{Wilson1974}
K.~G.~Wilson,
``Confinement of quarks,''
\textit{Phys.\ Rev.\ D}\ \textbf{10}, 2445 (1974).

\bibitem{Weisstein}
E.~W.~Weisstein,
``Almost integer,''
\textit{MathWorld}---A Wolfram Web Resource.

\bibitem{mpmath}
F.~Johansson et al.,
``mpmath: a Python library for arbitrary-precision floating-point arithmetic (version 1.3),''
\url{https://mpmath.org/} (2023).

\bibitem{Dyson1952}
F.~J.~Dyson,
``Divergence of perturbation theory in quantum electrodynamics,''
\textit{Phys.\ Rev.}\ \textbf{85}, 631 (1952).

\end{thebibliography}

\end{document}
